% --- Archon physics formalization source begin ---
% archon:physics
% archon:covers IPhO2026Problems/problem_IPhO_2026_3_C_3.lean
% archon:source-report reports/ipho_2026/problem_IPhO_2026_3_C_3.source.json
% archon:problem-id IPhO_2026_3
% archon:part-id C.3
% archon:previous-part IPhO_2026_3_B_1
% archon:previous-part-policy natural_language_prerequisite_only
% archon:previous-part IPhO_2026_3_C_2
% archon:previous-part-policy natural_language_prerequisite_only

\chapter{Physics problem IPhO\_2026\_3\_C\_3}
\label{ch:IPhO2026Problems_problem_IPhO_2026_3_C_3}

\paragraph{Problem source.}
The paramagnetic torus executes the Carnot refrigeration cycle
1 -> 2 -> 3 -> 4 -> 1 shown in Figure 3b in the H-versus-T plane.  T\_h and T\_c
are the hot- and cold-reservoir temperatures; Q\_h is the magnitude of heat
delivered to the hot reservoir and Q\_c is the magnitude absorbed from the cold
reservoir.  The equation of state is T*M*V = n*K*H and the isothermal heat
relation from part B may be reused.

Current subquestion:
Using the supplied potassium-chromate and liquid-helium data, find the helium temperature after one cycle.

\paragraph{Current subquestion.}
Using the supplied potassium-chromate and liquid-helium data, find the helium temperature after one cycle.

\paragraph{Recorded answer/context.}
Q\_c = 1.29e-1 J, so |Delta T| = 9.92e-3 K and T\_final = 0.99008 K.

\paragraph{Figure/image path.}
/root/proposal\_for\_physic/science-mango/ipho\_2026\_source/image/T3\_page-4.png

\paragraph{Reusable previous-part conclusions.}
\begin{itemize}
\item Source B.1. Question: At fixed temperature T, H changes from H\_i to H\_f. Find the heat Q transferred into the torus. Reusable conclusions: Q = -(mu\_0*n*K/(2*T))*(H\_f\textasciicircum{}2 - H\_i\textasciicircum{}2). Policy: natural\_language\_prerequisite\_only; do\_not\_import\_Lean\_output
\item Source C.2. Question: Express M\_1 in terms of M\_2, M\_3, and M\_4. Reusable conclusions: M\_1 = sqrt(M\_2\textasciicircum{}2 - M\_3\textasciicircum{}2 + M\_4\textasciicircum{}2), taking the nonnegative magnitude. Policy: natural\_language\_prerequisite\_only; do\_not\_import\_Lean\_output
\end{itemize}

\paragraph{Formalization target.}
Redraft the existing scaffold with an explicit Mathlib import in addition to its
Physlib units import.  Preserve the dimensioned thermodynamic and magnetic
quantities, Figure 3b states, supplied SI readouts, licensed B.1/C.2 relations,
calorimetry and numerical tolerances.  Leave proof bodies as `by sorry`.

\begin{definition}[CarnotState]
  \label{decl:physics:IPhO_2026_3_C_3:CarnotState}
  \lean{IPhO2026Problems.IPhO2026_3_C_3.CarnotState}
  The four vertices of the Carnot cycle “1 → 2 → 3 → 4 → 1”.
\end{definition}
\begin{proof}
  This is the defining declaration; unfolding it gives the stated typed data or relation.
\end{proof}

\begin{definition}[CarnotState.next]
  \label{decl:physics:IPhO_2026_3_C_3:CarnotState:next}
  \lean{IPhO2026Problems.IPhO2026_3_C_3.CarnotState.next}
  \uses{decl:physics:IPhO_2026_3_C_3:CarnotState}
  The oriented cycle order shown in Figure 3b.
\end{definition}
\begin{proof}
  This is the defining declaration; unfolding it gives the stated typed data or relation.
\end{proof}

\begin{definition}[Temperature]
  \label{decl:physics:IPhO_2026_3_C_3:Temperature}
  \lean{IPhO2026Problems.IPhO2026_3_C_3.Temperature}
  \uses{decl:physics:IPhO_2026_3_C_3:Volume}
  Absolute temperature, recorded in kelvin in the supplied-data hypotheses.
\end{definition}
\begin{proof}
  This is the defining declaration; unfolding it gives the stated typed data or relation.
\end{proof}

\begin{definition}[Volume]
  \label{decl:physics:IPhO_2026_3_C_3:Volume}
  \lean{IPhO2026Problems.IPhO2026_3_C_3.Volume}
  Volume, recorded in cubic metres in the supplied-data hypotheses.
\end{definition}
\begin{proof}
  This is the defining declaration; unfolding it gives the stated typed data or relation.
\end{proof}

\begin{definition}[MassDensity]
  \label{decl:physics:IPhO_2026_3_C_3:MassDensity}
  \lean{IPhO2026Problems.IPhO2026_3_C_3.MassDensity}
  Mass density, recorded in kilograms per cubic metre.
\end{definition}
\begin{proof}
  This is the defining declaration; unfolding it gives the stated typed data or relation.
\end{proof}

\begin{definition}[SpecificHeatCapacity]
  \label{decl:physics:IPhO_2026_3_C_3:SpecificHeatCapacity}
  \lean{IPhO2026Problems.IPhO2026_3_C_3.SpecificHeatCapacity}
  Specific heat capacity, with dimension “L² T⁻² Θ⁻¹”.
\end{definition}
\begin{proof}
  This is the defining declaration; unfolding it gives the stated typed data or relation.
\end{proof}

\begin{definition}[MagneticFieldStrength]
  \label{decl:physics:IPhO_2026_3_C_3:MagneticFieldStrength}
  \lean{IPhO2026Problems.IPhO2026_3_C_3.MagneticFieldStrength}
  \uses{decl:physics:IPhO_2026_3_C_3:Magnetization}
  Magnetic field strength “H”, with SI unit ampere per metre.
\end{definition}
\begin{proof}
  This is the defining declaration; unfolding it gives the stated typed data or relation.
\end{proof}

\begin{definition}[Magnetization]
  \label{decl:physics:IPhO_2026_3_C_3:Magnetization}
  \lean{IPhO2026Problems.IPhO2026_3_C_3.Magnetization}
  \uses{decl:physics:IPhO_2026_3_C_3:Energy}
  Magnetization “M”, with the same SI dimension as magnetic field strength.
\end{definition}
\begin{proof}
  This is the defining declaration; unfolding it gives the stated typed data or relation.
\end{proof}

\begin{definition}[Energy]
  \label{decl:physics:IPhO_2026_3_C_3:Energy}
  \lean{IPhO2026Problems.IPhO2026_3_C_3.Energy}
  Energy, recorded in joules.
\end{definition}
\begin{proof}
  This is the defining declaration; unfolding it gives the stated typed data or relation.
\end{proof}

\begin{definition}[MagneticPermeability]
  \label{decl:physics:IPhO_2026_3_C_3:MagneticPermeability}
  \lean{IPhO2026Problems.IPhO2026_3_C_3.MagneticPermeability}
  Magnetic permeability, with SI unit newton per ampere squared.
\end{definition}
\begin{proof}
  This is the defining declaration; unfolding it gives the stated typed data or relation.
\end{proof}

\begin{definition}[Setup]
  \label{decl:physics:IPhO_2026_3_C_3:Setup}
  \lean{IPhO2026Problems.IPhO2026_3_C_3.Setup}
  \uses{decl:physics:IPhO_2026_3_C_3:CarnotState, decl:physics:IPhO_2026_3_C_3:Temperature, decl:physics:IPhO_2026_3_C_3:Volume, decl:physics:IPhO_2026_3_C_3:MassDensity, decl:physics:IPhO_2026_3_C_3:SpecificHeatCapacity, decl:physics:IPhO_2026_3_C_3:MagneticFieldStrength, decl:physics:IPhO_2026_3_C_3:Magnetization, decl:physics:IPhO_2026_3_C_3:Energy, decl:physics:IPhO_2026_3_C_3:MagneticPermeability}
  All physical quantities needed to describe the torus, cycle, and helium sample.
\end{definition}
\begin{proof}
  This is the defining declaration; unfolding it gives the stated typed data or relation.
\end{proof}

\begin{definition}[HasSuppliedData]
  \label{decl:physics:IPhO_2026_3_C_3:HasSuppliedData}
  \lean{IPhO2026Problems.IPhO2026_3_C_3.HasSuppliedData}
  \uses{decl:physics:IPhO_2026_3_C_3:CarnotState, decl:physics:IPhO_2026_3_C_3:Setup}
  The numerical readouts supplied on the official source page, all expressed in the fixed SI units documented by the corresponding fields.
\end{definition}
\begin{proof}
  This is the defining declaration; unfolding it gives the stated typed data or relation.
\end{proof}

\begin{definition}[GoverningLaws]
  \label{decl:physics:IPhO_2026_3_C_3:GoverningLaws}
  \lean{IPhO2026Problems.IPhO2026_3_C_3.GoverningLaws}
  \uses{decl:physics:IPhO_2026_3_C_3:CarnotState, decl:physics:IPhO_2026_3_C_3:Setup}
  The governing physical model: Figure 3b's isotherm labels, material geometry, the paramagnet equation of state, and calorimetric energy conservation for the helium.
\end{definition}
\begin{proof}
  This is the defining declaration; unfolding it gives the stated typed data or relation.
\end{proof}

\begin{definition}[PreviousPartResults]
  \label{decl:physics:IPhO_2026_3_C_3:PreviousPartResults}
  \lean{IPhO2026Problems.IPhO2026_3_C_3.PreviousPartResults}
  \uses{decl:physics:IPhO_2026_3_C_3:CarnotState, decl:physics:IPhO_2026_3_C_3:Setup}
  The two reusable results explicitly licensed by the blueprint: the part B.1 isothermal-heat relation on the cold leg “2 → 3”, and the nonnegative-magnitude relation from part C.2.
\end{definition}
\begin{proof}
  This is the defining declaration; unfolding it gives the stated typed data or relation.
\end{proof}

\begin{theorem}[Physics formalization target]
\label{thm:physics:IPhO_2026_3_C_3:target}
\lean{IPhO2026Problems.IPhO2026_3_C_3.helium_temperature_after_one_cycle}
\uses{decl:physics:IPhO_2026_3_C_3:CarnotState, decl:physics:IPhO_2026_3_C_3:CarnotState:next, decl:physics:IPhO_2026_3_C_3:Temperature, decl:physics:IPhO_2026_3_C_3:Volume, decl:physics:IPhO_2026_3_C_3:MassDensity, decl:physics:IPhO_2026_3_C_3:SpecificHeatCapacity, decl:physics:IPhO_2026_3_C_3:MagneticFieldStrength, decl:physics:IPhO_2026_3_C_3:Magnetization, decl:physics:IPhO_2026_3_C_3:Energy, decl:physics:IPhO_2026_3_C_3:MagneticPermeability, decl:physics:IPhO_2026_3_C_3:Setup, decl:physics:IPhO_2026_3_C_3:HasSuppliedData, decl:physics:IPhO_2026_3_C_3:GoverningLaws, decl:physics:IPhO_2026_3_C_3:PreviousPartResults}
The supplied cycle data imply \(Q_c\approx0.129\) J, a helium temperature drop
of approximately \(0.00992\) K, and final temperature approximately
\(0.99008\) K, each within the stated tolerance.
\end{theorem}
\begin{proof}
Insert \(H_2,H_3,\mu_0,n,K,T_c\) into the licensed cold-isotherm heat formula
to obtain \(Q_c\).  The helium mass is its density times \(1.00\) L, so
calorimetry gives
\(\Delta T=Q_c/(\rho_{\mathrm{He}}V_{\mathrm{He}}c_{\mathrm{He}})\).
Subtract this positive drop from \(1.00\) K and check the three rounding
intervals.
\end{proof}
% --- Archon physics formalization source end ---
